\documentclass[11pt]{amsart}
\usepackage{amsmath,amssymb,amsthm}
\usepackage[margin=1.05in]{geometry}

\numberwithin{equation}{section}
\newtheorem{theorem}{Theorem}[section]
\newtheorem{proposition}[theorem]{Proposition}
\newtheorem{lemma}[theorem]{Lemma}
\newtheorem{corollary}[theorem]{Corollary}
\theoremstyle{definition}
\newtheorem{definition}[theorem]{Definition}
\theoremstyle{remark}
\newtheorem{remark}[theorem]{Remark}
\newtheorem{example}[theorem]{Example}

\newcommand{\Z}{\mathbb{Z}}
\newcommand{\Q}{\mathbb{Q}}
\newcommand{\R}{\mathbb{R}}
\newcommand{\N}{\mathbb{N}}
\newcommand{\C}{\mathbb{C}}
\newcommand{\T}{\mathbb{T}}
\newcommand{\cA}{\mathcal{A}}
\newcommand{\cC}{\mathcal{C}}
\newcommand{\cO}{\mathcal{O}}
\newcommand{\Ent}{\mathrm{S}}
\newcommand{\cR}{\mathcal{R}}
\newcommand{\pp}{\mathfrak{p}}
\newcommand{\qq}{\mathfrak{q}}
\newcommand{\aaa}{\mathfrak{a}}
\newcommand{\bb}{\mathfrak{b}}
\newcommand{\mm}{\mathfrak{m}}
\newcommand{\hZ}{\widehat{\mathbb{Z}}}
\newcommand{\Ns}{\mathrm{N}}
\newcommand{\Cl}{\operatorname{Cl}}
\newcommand{\Frob}{\operatorname{Frob}}
\newcommand{\Gal}{\operatorname{Gal}}
\newcommand{\Prob}{\operatorname{Prob}}
\newcommand{\Xiext}{\Xi^{\mathrm{ext}}_\beta}
\newcommand{\nusym}{\nu_{\mathrm{sym}}}
\newcommand{\dw}{\mathfrak{d}}

\PassOptionsToPackage{hyphens}{url}
\usepackage[hidelinks,bookmarks=false]{hyperref}
\hypersetup{pdftitle={Localized Bost-Connes systems I: factor type and the spectrum of transitions},pdfauthor={Ruqing Chen}}

\title[Factor type and the spectrum of transitions]{Localized Bost--Connes systems I:\\ factor type and the spectrum of transitions}
\author{Ruqing Chen}
\address{GUT Geoservice Inc., Montreal, Canada}
\email{ruqing@hotmail.com}
\thanks{DOI: \href{https://doi.org/10.5281/zenodo.22160827}{10.5281/zenodo.22160827}.
Sources, code and data:\newline\url{https://github.com/Ruqing1963/localized-bost-connes}
(this paper: \texttt{papers/I-type-spectrum/}).}
\date{}

\begin{document}

\begin{abstract}
The obstruction group $\Xi_\beta(K,S)$ of \cite{Main} classifies the $\mathrm{KMS}_\beta$
simplex. We show that it is the $s=0$ slice of a single larger object, the \emph{extended
obstruction group}
\[
\Xiext(K,S)=\Bigl\{(\chi,s)\in\widehat{G_S}\times\R:\ \sum_{\pp\in S}\Ns\pp^{-\beta}\bigl|1-\chi(\sigma_\pp)\Ns\pp^{\,i\beta s}\bigr|^2<\infty\Bigr\},
\]
whose slices $s\ne0$ control the Murray--von Neumann--Connes type of the associated
factors from above: the essential values of the Radon--Nikodym cocycle annihilate
$\Sigma_\beta=\mathrm{pr}_\R\Xiext$, so $\Sigma_\beta\ne\{0\}$ excludes type
$\mathrm{III}_1$, a dense $\Sigma_\beta$ forces $\mathrm{III}_0$, and $\Sigma_\beta=c\Z$
forces $\mathrm{III}_{\lambda^k}$ with $\lambda=e^{-2\pi/c}$ or $\mathrm{III}_0$; moreover
$\Sigma_\beta$ lies in the point spectrum of the flow of weights. Matching lower bounds come
from a direct construction of essential values in the manner of Araki and Woods. The phase
classification and the type classification are thus two coordinate slices of one
Dirichlet-series condition. For $S=P_K$ and $\beta=1$ both reduce to the non-vanishing of Hecke $L$-functions
on $\operatorname{Re}=1$: the $s=0$ slice is $L(1,\chi)\ne0$ and the $s\ne0$ slices are
$L(1-i\tau,\chi)\ne0$.

Two consequences are perhaps unexpected. First, the natural sufficient condition for type
$\mathrm{III}_1$ --- that two primes have irrationally related log-norms --- is false: two
primes contribute a finite amount to a divergent series. We exhibit $S\subset P_\Q$ with
$\zeta_S(1)=\infty$, a \emph{unique} $\mathrm{KMS}_1$ state, and factors of pure type
$\mathrm{III}_{e^{-1}}$; uniqueness of the KMS state and type $\mathrm{III}_1$ are logically
independent. Second, for infinite $S$ the log-norms are never pairwise commensurable, so any
failure of $\mathrm{III}_1$ comes from \emph{approximate} commensurability along a set of
divergent $\beta$-weight.

In the horizontal direction we determine which chains $\{\Xi_\beta\}_{\beta>0}$ occur. For
every sequence $(\beta_n)\subset(0,1)$ there is $S$ with $\beta_c(S)=1$, dense parameter
subgroup, and $\Xi_\beta=\bigoplus_{n:\beta_n<\beta}\Z/2\Z$; taking $(\beta_n)$ dense in
$(0,1)$ gives countably many dense transitions, and taking it dense in a Cantor set $\cC$
makes the non-local-constancy locus exactly $\cC$. The continuous devil's staircase, however,
is unattainable for a structural reason: $\widehat{G_S}$ is countable for every $S$, so the
jump set is countable and every order parameter is a pure jump function. It is approachable
in the uniform norm and not attained.
\end{abstract}

\maketitle

\tableofcontents

\section{Introduction: one group, two directions}

We use the notation of \cite{Main}: $\cA_{K,S}=C(Y_S)\rtimes J_S$ with
$\sigma_t(\mu_\aaa)=\Ns\aaa^{it}\mu_\aaa$, symmetry group $G_S$, parameter homomorphism
$\sigma:I_S\to G_S$, and
\begin{equation}\label{eq:Xi}
\Xi_\beta(K,S)=\Bigl\{\chi\in\widehat{G_S}:\ \sum_{\pp\in S}\Ns\pp^{-\beta}\bigl|1-\chi(\sigma_\pp)\bigr|^2<\infty\Bigr\},
\end{equation}
the $\mathrm{KMS}_\beta$ simplex being $\Prob\bigl(G_S/\Xi_\beta^\perp\bigr)$. Write
$\cR_S$ for the tail equivalence relation on $V=\Z_{\ge0}^S$ and $\pi_\beta$ for the forced
product measure \cite[Lem.~2.10]{Main}.

\emph{The series.} The foundational paper is \cite{Main}. The present paper is the first
of a series continuing it; the later papers, in preparation, treat the ideal structure and
boundary quotient, the invariants that see $\Xi_\beta$ and those that do not, the metric
geometry, and the free variant. Nothing here depends on them.

The phase diagram of $\cA_{K,S}$ has two directions. Horizontally one asks which
$\beta$ are transition points, i.e.\ how the chain $\{\Xi_\beta\}$ grows; vertically one asks
what the factors at a given $\beta$ are. Definition~\ref{def:Xiext} shows the two questions
are slices of one object, and the rest of the paper answers both --- the second one
partly: the object bounds the type from above, and the bound is matched by direct
constructions in the cases treated.

\begin{definition}\label{def:Xiext}
The \emph{extended obstruction group} is
\begin{equation}\label{eq:Xiext}
\Xiext(K,S)=\Bigl\{(\chi,s)\in\widehat{G_S}\times\R:\
\sum_{\pp\in S}\Ns\pp^{-\beta}\bigl|1-\chi(\sigma_\pp)\,\Ns\pp^{\,i\beta s}\bigr|^2<\infty\Bigr\},
\end{equation}
a subgroup of $\widehat{G_S}\times\R$, and $\Sigma_\beta:=\mathrm{pr}_\R(\Xiext)\le\R$.
\end{definition}

That $\Xiext$ is a subgroup follows as in \cite[Lem.~3.5]{Main} from
$|1-z_1z_2|\le|1-z_1|+|1-z_2|$ for $|z_i|=1$; and $\Xi_\beta=\Xiext\cap(\widehat{G_S}\times\{0\})$.

\section{The type criterion}\label{sec:type}

Let $\beta>0$ with $\zeta_{K,S}(\beta)=\infty$, let $\varphi$ be an extremal
$\mathrm{KMS}_\beta$ state and $M=\pi_\varphi(\cA_{K,S})''$. By \cite[Prop.~2.8 and
Lem.~2.10]{Main}, $\varphi$ is given by a probability measure $\nu$ on $Y=V\times G_S$ whose
$V$-marginal is $\pi_\beta$, and $M$ is the von Neumann algebra of the orbit equivalence
relation $\cR$ on $(Y,\nu)$ generated by the maps $(v,h)\mapsto(v+e_\pp,\sigma_\pp h)$,
$\pp\in S$. The Radon--Nikodym cocycle of $\cR$ is
\begin{equation}\label{eq:rho}
\rho\bigl(y,\aaa\bb^{-1}\cdot y\bigr)=\log\frac{d\nu\circ(\aaa\bb^{-1})}{d\nu}(y)
=-\beta\log\frac{\Ns\aaa}{\Ns\bb},
\end{equation}
and we write $E\subseteq\R$ for its group of finite essential values in the sense of
\cite{Schmidt}: $u\in E$ iff for every $B\subseteq Y$ with $\nu(B)>0$ and every
$\varepsilon>0$ there is a partial transformation $T$ of $\cR$ with
$\nu(B\cap T^{-1}B)>0$ and $|\rho(y,Ty)-u|<\varepsilon$ on its domain. $E$ is a closed
subgroup of $\R$ \cite{Schmidt}, $M$ is a factor of type $\mathrm{III}$
\cite[Prop.~6.1]{Main}, and by Krieger's theorem \cite{Krieger} the Connes invariant is
\begin{equation}\label{eq:SM}
S(M)=\{0\}\cup\exp(E),
\end{equation}
so that $M$ is of type $\mathrm{III}_1$, $\mathrm{III}_\lambda$ or $\mathrm{III}_0$
according as $E=\R$, $E=(\log\lambda)\Z$ with $0<\lambda<1$, or $E=\{0\}$ \cite{Connes73}.
Everything in this section is a computation of $E$.

\begin{lemma}[Coboundary dictionary]\label{lem:dict}
Let $(\chi,s)\in\widehat{G_S}\times\R$ and $c_\pp=\chi(\sigma_\pp)\Ns\pp^{\,i\beta s}$. The
following are equivalent:
\begin{enumerate}
\item there is a measurable $g:V\to\T$ with $g(v+e_\pp)=c_\pp\,g(v)$ for $\pi_\beta$-a.e.\
$v$, for every $\pp\in S$;
\item $(\chi,s)\in\Xiext$.
\end{enumerate}
\end{lemma}

\begin{proof}
Propositions 3.3 and 3.4 of \cite{Main} prove exactly this equivalence for
$c_\pp=\overline{\chi(\sigma_\pp)}$, and their proofs use only that $|c_\pp|=1$, that the
coordinates $v_\pp$ are independent under $\pi_\beta$ with $\pi_\beta(v_\pp=k)=(1-t_\pp)t_\pp^k$,
$t_\pp=\Ns\pp^{-\beta}$, and the estimate $1-|m_\pp|^2\asymp t_\pp|1-c_\pp|^2$ of
\cite[Lem.~3.2]{Main} for $m_\pp=(1-t_\pp)/(1-c_\pp t_\pp)$, which holds for every unimodular
$c_\pp$. (Conjugating $g$ replaces $c_\pp$ by $\overline{c_\pp}$, and
$|1-c_\pp|=|1-\overline{c_\pp}|$.)
\end{proof}

\begin{proposition}[The extended group bounds the type]\label{prop:bound}
Put $\Sigma_\beta^\vee=\{u\in\R:\ su\in2\pi\Z\ \text{for all}\ s\in\Sigma_\beta\}$. Then
$E\subseteq\Sigma_\beta^\vee$. Consequently:
\begin{enumerate}
\item if $\Sigma_\beta\ne\{0\}$, then $M$ is not of type $\mathrm{III}_1$;
\item if $\Sigma_\beta$ is dense in $\R$, then $M$ is of type $\mathrm{III}_0$;
\item if $\Sigma_\beta=c\Z$ with $c>0$, then $M$ is of type $\mathrm{III}_{\lambda^k}$ for
some integer $k\ge1$, where $\lambda=e^{-2\pi/c}$, or of type $\mathrm{III}_0$.
\end{enumerate}
\end{proposition}

\begin{proof}
Let $(\chi,s)\in\Xiext$, let $g$ be as in Lemma~\ref{lem:dict}, and put
$G(v,h)=g(v)\overline{\chi(h)}$, a unimodular measurable function on $Y$. For $\aaa\in J_S$,
\[
G(v+a,\sigma_\aaa h)=\chi(\sigma_\aaa)\Ns\aaa^{\,i\beta s}g(v)\,\overline{\chi(\sigma_\aaa)}\,\overline{\chi(h)}
=\Ns\aaa^{\,i\beta s}\,G(v,h),
\]
so by \eqref{eq:rho} the $\T$-valued cocycle $e^{-is\rho}$ on $\cR$ is the coboundary
$e^{-is\rho(y,y')}=G(y')/G(y)$. A coboundary has no essential value other than $1$: if
$|G-\gamma|<\delta$ on a set $B_\delta$ of positive measure, every pair $y,y'\in B_\delta$
has $|G(y')/G(y)-1|<2\delta$. On the other hand, if $u\in E$ then $e^{-isu}$ is an essential
value of $e^{-is\rho}$, since $|e^{-is\rho}-e^{-isu}|\le|s|\,|\rho-u|$. Hence
$e^{-isu}=1$ for all $u\in E$ and all $s\in\Sigma_\beta$, i.e.\ $E\subseteq\Sigma_\beta^\vee$.

(1) If $\Sigma_\beta\ne\{0\}$ then $\Sigma_\beta^\vee\ne\R$, so $E\ne\R$. (2) If $\Sigma_\beta$
is dense then $\Sigma_\beta^\vee=\{0\}$, so $E=\{0\}$ and $S(M)=\{0,1\}$. (3) If
$\Sigma_\beta=c\Z$ then $\Sigma_\beta^\vee=\frac{2\pi}{c}\Z$, and the closed subgroups of
$\frac{2\pi}{c}\Z$ are $\{0\}$ and $\frac{2\pi k}{c}\Z$, $k\ge1$; by \eqref{eq:SM} these give
$S(M)=\{0,1\}$ and $S(M)=\{0\}\cup(e^{-2\pi k/c})^{\Z}$ respectively.
\end{proof}

\begin{remark}[$\Sigma_\beta$ and the flow of weights]\label{rem:flow}
The proof shows more. With $G$ as above, $\Psi(y,u)=G(y)e^{isu}$ is invariant under the
Maharam extension $(y,u)\mapsto(y',u+\rho(y,y'))$ of $\cR$ \cite{Maharam} and satisfies
$\Psi(y,u+t)=e^{ist}\Psi(y,u)$; so every $s\in\Sigma_\beta$ is an eigenvalue of the flow of
weights of $M$, i.e.\ $\Sigma_\beta$ is contained in the point spectrum of that flow, which by
Connes--Takesaki \cite{CT} is Connes' invariant $T(M)$. We shall not need the converse
inclusion and do not prove it; it holds when the flow of weights has pure point spectrum,
and in that case $\Sigma_\beta$ alone determines the type. Whether the flow of weights of
$M$ always has pure point spectrum is not known to us, and this is why
Proposition~\ref{prop:bound} is only an upper bound. Lower bounds are obtained by exhibiting
essential values directly, which is the purpose of the next lemma.
\end{remark}

\begin{lemma}[Asymptotic pairs]\label{lem:pairs}
Suppose $\Xi_\beta=\{1\}$, so that $\nu=\pi_\beta\otimes\mathrm{Haar}$ \cite[Thm.~3.6]{Main}.
Let $u\in\R$ and let $(\aaa_j,\bb_j)_{j\ge1}$ be pairs of primes of $S$, all the primes
$\aaa_j,\bb_j$ distinct, such that
\[
\beta\log\frac{\Ns\aaa_j}{\Ns\bb_j}\to u,\qquad
g_j:=\sigma_{\bb_j}\sigma_{\aaa_j}^{-1}\to1\ \text{in}\ G_S,\qquad
\sum_j\Ns\aaa_j^{-\beta}=\infty .
\]
Then $u\in E$.
\end{lemma}

\begin{proof}
Write $t_\pp=\Ns\pp^{-\beta}$ and, for each $j$,
$P_j=\{v_{\aaa_j}=1,\,v_{\bb_j}=0\}$, $Q_j=\{v_{\aaa_j}=0,\,v_{\bb_j}=1\}$, events of
$\pi_\beta$-probability $p_j=(1-t_{\aaa_j})(1-t_{\bb_j})t_{\aaa_j}$ and
$q_j=(1-t_{\aaa_j})(1-t_{\bb_j})t_{\bb_j}$. The events $P_j\cup Q_j$ are independent and
$\sum_jp_j=\infty$; since $t_{\bb_j}/t_{\aaa_j}\to e^{u}$, there is $\theta>0$ with
$p_j\ge\theta(p_j+q_j)$ for all large $j$. Fix $\varepsilon>0$ and $j_0$, and let $D$ be the
set of $(v,h)$ such that some $j\ge j_0$ has $v\in P_j\cup Q_j$ and the least such $j$,
denoted $j(v)$, has $v\in P_j$. Define
\[
T(v,h)=\bigl(v-e_{\aaa_j}+e_{\bb_j},\ g_jh\bigr),\qquad j=j(v),
\]
on $D$. The move changes only the coordinates $\aaa_j,\bb_j$, turning $P_j$ into $Q_j$, so
$T$ is injective with image the set where the least $j\ge j_0$ with $v\in P_j\cup Q_j$ has
$v\in Q_j$; $T$ is a partial transformation of $\cR$, being the map $\bb_j\aaa_j^{-1}$ on
each piece $D_j=\{j(v)=j\}$, and by \eqref{eq:rho}, $\rho(y,Ty)=\beta\log(\Ns\aaa_j/\Ns\bb_j)$
lies in $(u-\varepsilon,u+\varepsilon)$ once $j_0$ is large. Since the $D_j$ are disjoint
and $\sum_j(p_j+q_j)=\infty$,
\[
\nu(D)=\sum_{j\ge j_0}p_j\prod_{j_0\le i<j}(1-p_i-q_i)
\ \ge\ \theta\sum_{j\ge j_0}(p_j+q_j)\prod_{j_0\le i<j}(1-p_i-q_i)=\theta .
\]

Now let $\nu(B)>0$ and $0<\eta<1$. The sets measurable with respect to finitely many
coordinates of $V$ together with the $G_S$-coordinate generate the $\sigma$-algebra of $Y$,
so there are a finite $F\subset S$ and $A\in\sigma(v_\pp:\pp\in F)\otimes\mathcal B(G_S)$ with
$\nu(A\,\triangle\,B)<\eta\,\nu(B)$. Translation in the second coordinate preserves $\nu$
and is continuous in $L^1(\nu)$, so $\nu(A\,\triangle\,(1\times g)^{-1}A)\to0$ as $g\to1$.
Choose $j_0$ so large that $\aaa_j,\bb_j\notin F$, $|\rho-u|<\varepsilon$ on $D_j$, and
$\nu(A\,\triangle\,(1\times g_j)^{-1}A)<\eta\,\nu(B)$ for all $j\ge j_0$. For
$(v,h)\in D_j$ the point $T(v,h)$ agrees with $(v,g_jh)$ on the coordinates in $F$, hence
$D_j\cap A\cap(1\times g_j)^{-1}A\subseteq A\cap T^{-1}A$; the set $D_j$ depends only on
the coordinates $\aaa_i,\bb_i$, $i\ge j_0$, so it is independent of
$A\cap(1\times g_j)^{-1}A$, and
\[
\nu(A\cap T^{-1}A)\ \ge\ \sum_{j\ge j_0}\nu(D_j)\,\nu\bigl(A\cap(1\times g_j)^{-1}A\bigr)
\ \ge\ \nu(D)\,\bigl(\nu(A)-\eta\nu(B)\bigr)\ \ge\ \theta(1-2\eta)\,\nu(B).
\]
Finally $d(\nu\circ T^{-1})/d\nu=e^{-\rho\circ T^{-1}}\le e^{|u|+\varepsilon}$ on $T(D)$, so
\[
\nu(B\cap T^{-1}B)\ \ge\ \nu(A\cap T^{-1}A)-\nu(A\triangle B)-\nu\bigl(T^{-1}(A\triangle B)\bigr)
\ \ge\ \bigl(\theta(1-2\eta)-\eta-e^{|u|+\varepsilon}\eta\bigr)\nu(B)>0
\]
for $\eta$ small. As $B$ and $\varepsilon$ were arbitrary, $u\in E$.
\end{proof}

The lemma is the asymptotic-ratio-set argument of Araki and Woods \cite{AW} with the
$G_S$-coordinate carried along; the hypothesis $g_j\to1$ is what makes that coordinate
harmless. The restriction to $\Xi_\beta=\{1\}$ is one of convenience: it is the case needed
below, and it is the case in which $\nu$ has the product form used in the last display.

\section{The classical systems}\label{sec:classical}

\begin{proposition}[The slices for $S=P_K$]\label{prop:slices}
Let $S=P_K$ and $0<\beta\le1$. Then $\Xiext(K,S)=\{(1,0)\}$; in particular
$\Sigma_\beta=\{0\}$.
\end{proposition}

\begin{proof}
For $(\chi,s)$ put $\tau=\beta s$; since $|1-z|^2=2-2\operatorname{Re}z$ for $|z|=1$,
\begin{equation}\label{eq:twosums}
\sum_{\Ns\pp\le x}\Ns\pp^{-\beta}\bigl|1-\chi(\sigma_\pp)\Ns\pp^{\,i\tau}\bigr|^2
=2\sum_{\Ns\pp\le x}\Ns\pp^{-\beta}-2\operatorname{Re}\sum_{\Ns\pp\le x}\chi(\sigma_\pp)\Ns\pp^{-\beta+i\tau},
\end{equation}
and it suffices to show that for $(\chi,\tau)\ne(1,0)$ the second sum is $o$ of the first.
A character of $G_{P_K}=\Gal(K^{ab}/K)$ is a ray class character, and by
\cite[Lem.~3.8]{Main} $\chi(\sigma_\pp)=\overline{\chi(\Frob_\pp)}$ off the conductor. For
$\beta=1$ the series $\sum_\pp\chi(\sigma_\pp)\Ns\pp^{-1+i\tau}$ differs from
$\log L(1-i\tau,\bar\chi)$ by an absolutely convergent sum over prime powers, and
$L(1-i\tau,\bar\chi)$ is finite --- the only pole of a Hecke $L$-function on
$\operatorname{Re}=1$ being the simple pole of $\zeta_K$ at $s=1$, excluded --- and nonzero
by the theorem of Hecke and Landau \cite{Hecke,Landau}. So the second sum is $O(1)$ while
$\sum_{\Ns\pp\le x}\Ns\pp^{-1}\to\infty$. For $0<\beta<1$, partial summation from the prime
ideal theorem in ray classes in the standard zero-free region \cite{Landau} gives
$O_A\bigl(x^{1-\beta}\log^{-A}x\bigr)$ for the second sum against
$\sum_{\Ns\pp\le x}\Ns\pp^{-\beta}\sim x^{1-\beta}/((1-\beta)\log x)$ for the first.
\end{proof}

\begin{theorem}[The classical systems are of type $\mathrm{III}_1$]\label{thm:classical}
For $S=P_K$ and $0<\beta\le1$, $M$ is of type $\mathrm{III}_1$.
\end{theorem}

For $K=\Q$ this is due to Bost and Connes \cite{BC}, and in general to Laca, Larsen and
Neshveyev \cite{LLN}; we include a proof because it is the model for the computation of
\S\ref{sec:counter}.

\begin{proof}
$\Xi_\beta=\{1\}$ by \cite[Cor.~3.10]{Main}, so Lemma~\ref{lem:pairs} applies, and we must
produce, for each $u\in\R$, pairs $(\aaa_j,\bb_j)$ with $\beta\log(\Ns\aaa_j/\Ns\bb_j)\to u$,
$\sigma_{\bb_j}\sigma_{\aaa_j}^{-1}\to1$ in $\Gal(K^{ab}/K)$ and $\sum_j\Ns\aaa_j^{-\beta}=\infty$.
Fix moduli $\mm_1\mid\mm_2\mid\cdots$ whose ray class fields exhaust $K^{ab}$
\cite[Ch.~VI]{Neu}; for primes
$\aaa,\bb$ in the same ray class modulo $\mm_k$ the element
$\sigma_\bb\sigma_\aaa^{-1}=\Frob_\aaa\Frob_\bb^{-1}$ is trivial on the ray class field of
$\mm_k$, so pairs chosen in a common class modulo $\mm_{k}$ with $k\to\infty$ satisfy
$\sigma_{\bb_j}\sigma_{\aaa_j}^{-1}\to1$. By the prime ideal theorem in ray classes
\cite{Landau}, for fixed $\mm$ and each class $\mathfrak c$ the number of primes
$\pp\in\mathfrak c$ with $\Ns\pp\in(y,y(1+\delta)]$ is $\asymp_{\mm,\delta}y/\log y$ for
$y\ge y_0(\mm,\delta)$. Let $r=e^{u/\beta}$. Use only every $N_0$-th $n$, with $N_0$ so large that the ranges
$[2^n,2^{n+1})\cup r^{-1}[2^n,2^{n+1})$ of the blocks used are pairwise disjoint. Divide
$[2^n,2^{n+1})$ into intervals $I=(y,y(1+\delta_n)]$ with $\delta_n\downarrow0$ slowly, and
for each $I$ and each class
$\mathfrak c$ modulo $\mm_{k_n}$, with $k_n\to\infty$ slowly enough that
$2^n\ge y_0(\mm_{k_n},\delta_n)$, pair off primes of $\mathfrak c$ with norm in $I$ against
primes of $\mathfrak c$ with norm in $r^{-1}I$, injectively and disjointly from earlier
pairs; the pairs of different blocks are then disjoint. Each pair has
$|\beta\log(\Ns\aaa/\Ns\bb)-u|\le2\beta\delta_n$, and the pairs with
$\Ns\aaa\in[2^n,2^{n+1})$ number $\gg2^n/n$ (for large $n$, a positive proportion of the primes
in $[2^n,2^{n+1})$ is used), so they contribute $\gg2^{n(1-\beta)}/n\ge1/n$ to
$\sum_j\Ns\aaa_j^{-\beta}$, which diverges along the blocks used. Lemma~\ref{lem:pairs} gives $u\in E$
for every $u$, so $E=\R$ and $M$ is of type $\mathrm{III}_1$ by \eqref{eq:SM}.
\end{proof}

\begin{remark}[Two slices of one analytic input]\label{rem:unify}
At $\beta=1$ and $S=P_K$ the $s=0$ slice of $\Xiext$ is the uniqueness half of the classical
theorem, equivalent by \cite[Rem.~3.11]{Main} to $L(1,\chi)\ne0$; the slices $s\ne0$ are
governed, by Proposition~\ref{prop:slices}, by $L(1-i\tau,\chi)\ne0$, and their triviality
is the necessary condition for type $\mathrm{III}_1$ of Proposition~\ref{prop:bound}(1).
The classification of the phase and the type of the factor are thus the two coordinate
slices of one obstruction group, and for the full prime set both are governed by the
non-vanishing of Hecke $L$-functions on the line $\operatorname{Re}=1$. That the
$s\ne0$ slices are trivial is, by Remark~\ref{rem:flow}, the statement that the flow of
weights has no nonzero eigenvalue; that $M$ is $\mathrm{III}_1$, i.e.\ that the flow is
trivial, needed the additional input of Lemma~\ref{lem:pairs}.
\end{remark}

\section{Two incommensurable norms do not suffice}\label{sec:counter}

\begin{lemma}[No exact commensurability]\label{lem:noexact}
Let $S$ be infinite. Then the numbers $\log\Ns\pp$, $\pp\in S$, are never pairwise
commensurable.
\end{lemma}

\begin{proof}
$\Ns\pp=p^{f_\pp}$ with $f_\pp\le[K:\Q]$, and only finitely many primes of $K$ lie above a
given $p$, so an infinite $S$ meets infinitely many rational primes; and
$\log p_1/\log p_2\notin\Q$ for distinct rational primes since $p_1^m=p_2^n$ forces $m=n=0$.
\end{proof}

So exact commensurability, which would put all the values of the Radon--Nikodym cocycle
\eqref{eq:rho} in a lattice, never occurs for infinite $S$, and any failure of type
$\mathrm{III}_1$ must come from \emph{approximate} commensurability. It does.

\begin{theorem}[Exact determination of a $\mathrm{III}_{e^{-1}}$ system]\label{thm:counter}
Let $K=\Q$, $\beta=1$, $\|\cdot\|$ the distance to the nearest integer, and
\begin{equation}\label{eq:band}
S=\{3\}\cup\bigl\{\ell\ \text{prime}:\ \ell\ge16,\ \|\log\ell\|\le\varepsilon(\ell)\bigr\},
\qquad\varepsilon(\ell)=(\log\log\ell)^{-1/2}.
\end{equation}
Then $\zeta_S(1)=\infty$, $S$ contains primes with irrationally related logarithms, and
\[
\Xi^{\mathrm{ext}}_1(\Q,S)=\{1\}\times2\pi\Z,\qquad \Sigma_1=2\pi\Z .
\]
Hence $\Xi_1=\{1\}$ --- the $\mathrm{KMS}_1$ state is \emph{unique} --- while $M$ is of pure
type $\mathrm{III}_{e^{-1}}$; both $\mathrm{III}_1$ and $\mathrm{III}_0$ are excluded.
\end{theorem}

The proof occupies the rest of this section. The mechanism is a band decomposition: because
\eqref{eq:band} localizes $\log\ell$ near an \emph{integer}, the rounded logarithm is
constant on each band and the exponential sums reduce to geometric series.

\begin{lemma}[Bands]\label{lem:bands}
Put $\varepsilon_n=(\log n)^{-1/2}$ and $n_0=55$, so $\varepsilon_n<\tfrac12$ for $n>n_0$.
For $n>n_0$ let $B_n=\{\ell\in S:|\log\ell-n|\le\varepsilon(\ell)\}$ and
$S_0=\{\ell\in S:\ell\le e^{n_0+1}\}$. Then:
\begin{enumerate}
\item $S=S_0\sqcup\bigsqcup_{n>n_0}B_n$ and $\sum_{\ell\in S_0}\ell^{-1}=O(1)$;
\item for $\ell\in B_n$, $\log\ell=n+\delta_\ell$ with $|\delta_\ell|\le\varepsilon_n(1+O(n^{-1}))$,
so $\mathrm{round}(\log\ell)=n$ is constant on $B_n$;
\item $w_n:=\sum_{\ell\in B_n}\ell^{-1}=\frac{2\varepsilon_n}{n}(1+o(1))=\frac{2}{n\sqrt{\log n}}(1+o(1))$,
eventually decreasing to $0$ up to an absolutely summable perturbation;
\item $W(x):=\sum_{\ell\in S,\ \ell\le x}\ell^{-1}=4\sqrt{\log\log x}+O(1)\to\infty$.
\end{enumerate}
\end{lemma}

\begin{proof}
(1) For $n>n_0$ the bands lie in the pairwise disjoint intervals $[e^{n-1/2},e^{n+1/2}]$;
for $\ell\le e^{n_0+1}$ one has $\varepsilon(\ell)>1/2$ and the condition may hold for every
$\ell$, so $S_0$ can be all primes up to $e^{n_0+1}$, of Mertens sum $O(1)$.
(2) $|\log\ell-n|<\tfrac12$ gives the rounding, and
$\log\log\ell=\log n+O(\varepsilon_n/n)$ gives $\varepsilon(\ell)=\varepsilon_n(1+O(n^{-1}))$.
(3) By (2), $B_n$ is the set of primes in $[e^{n-\varepsilon_n^-},e^{n+\varepsilon_n^+}]$;
Mertens with error $O(e^{-c\sqrt{\log y}})$ gives
$w_n=\log\frac{n+\varepsilon_n^+}{n-\varepsilon_n^-}+O(e^{-c\sqrt n})$, which is the stated
asymptotic, with main term $2n^{-1}(\log n)^{-1/2}$ strictly decreasing and both errors
absolutely summable.
(4) $\sum_{n\le N}2n^{-1}(\log n)^{-1/2}=4\sqrt{\log N}+O(1)$ by comparison with the
integral; put $N=\lfloor\log x\rfloor$.
\end{proof}

\begin{lemma}[Linearization]\label{lem:lin}
Let $\chi$ factor through $(\Z/M\Z)^\times$ with Dirichlet character $\tilde\chi$, so that
$\chi(\sigma_\ell)=\tilde\chi(\ell)$ for $\ell\nmid M$ by \cite[Lem.~2.10(1)]{Main}, and let
$F$ be the set of primes dividing $M$. For $s\in\R$ and $N=\lfloor\log x\rfloor$,
\begin{equation}\label{eq:linear}
\sum_{\ell\in S,\ \ell\le x}\frac{\bigl|1-\chi(\sigma_\ell)\ell^{is}\bigr|^2}{\ell}
=2W(x)-2\operatorname{Re}\Lambda(x)+O_M(1),
\qquad
\Lambda(x)=\sum_{n\le N}e^{isn}c_n+O\Bigl(|s|\sum_{n\le N}\varepsilon_nw_n\Bigr),
\end{equation}
where $c_n=\sum_{\ell\in B_n\setminus F}\tilde\chi(\ell)\,\ell^{-1}$. Moreover
\begin{equation}\label{eq:err}
|s|\sum_{n\le N}\varepsilon_nw_n=2|s|\log\log N+O(|s|)=o\bigl(W(x)\bigr).
\end{equation}
\end{lemma}

\begin{proof}
$|1-z|^2=2-2\operatorname{Re}z$ for $|z|=1$ gives the first identity. For $\ell\in B_n$,
$\ell^{is}=e^{is(n+\delta_\ell)}=e^{isn}+O(|s|\varepsilon_n)$ by $|e^{i\theta}-1|\le|\theta|$,
and summing against $\ell^{-1}$ over $B_n$ gives the second. For \eqref{eq:err},
$\varepsilon_nw_n=2\varepsilon_n^2/n\,(1+o(1))=2/(n\log n)(1+o(1))$ and
$\sum_{n\le N}(n\log n)^{-1}=\log\log N+O(1)$; since $W(x)=4\sqrt{\log N}+O(1)$ and
$\log\log N=o(\sqrt{\log N})$, the claim follows. In the variable $x$ this reads
\begin{equation}\label{eq:absorb}
|s|\sum_{n\le N}\varepsilon_nw_n\ \asymp\ 2|s|\log\log\log x
\ =\ o\Bigl(\sqrt{\log\log x}\Bigr)\ =\ o\bigl(W(x)\bigr),
\end{equation}
the comparison being with $\sqrt{\log\log x}$ and not with $\log\log x$, since
$W(x)=4\sqrt{\log\log x}+O(1)$ by Lemma~\ref{lem:bands}(4).
\end{proof}

\begin{remark}[The one delicate point]\label{rem:delicate}
Estimate \eqref{eq:err} is where care is needed: the error from linearizing $\ell^{is}$
across a band \emph{diverges}, at rate $\log\log\log x$, and is absorbed only because $W(x)$
diverges faster, at rate $\sqrt{\log\log x}$. The exponent $-\tfrac12$ in
$\varepsilon(\ell)=(\log\log\ell)^{-1/2}$ is exactly what makes the convergence in
Proposition~\ref{prop:trivialchi} and this absorption hold simultaneously: an exponent
$-\tfrac13$ would break the former and $-1$ the latter.
\end{remark}

\begin{lemma}[Cancellation in bands]\label{lem:SW}
If $\tilde\chi$ is non-principal modulo $M$ then $|c_n|\ll_M e^{-c\sqrt n}$, so
$\sum_n|c_n|<\infty$.
\end{lemma}

\begin{proof}
Siegel--Walfisz gives
$\sum_{p\le y,\ p\equiv a\,(M)}p^{-1}=\varphi(M)^{-1}\log\log y+C(M,a)+O_M(e^{-c\sqrt{\log y}})$.
Applying this at the two endpoints of $B_n$, whose logarithms are $n+O(1)$, the constants
$C(M,a)$ cancel and
$\sum_{\ell\in B_n,\ \ell\equiv a}\ell^{-1}=w_n/\varphi(M)+O_M(e^{-c\sqrt n})$. Hence
$c_n=\frac{w_n}{\varphi(M)}\sum_a^{*}\tilde\chi(a)+O_M(\varphi(M)e^{-c\sqrt n})$,
and the character sum vanishes.
\end{proof}

\begin{proposition}[Non-trivial characters]\label{prop:chinontriv}
If $\chi\ne1$ then $\sum_{\ell\in S}\ell^{-1}|1-\chi(\sigma_\ell)\ell^{is}|^2=\infty$ for
\emph{every} $s\in\R$. In particular $\Xi^{\mathrm{ext}}_1\subseteq\{1\}\times\R$ and, taking
$s=0$, $\Xi_1=\{1\}$.
\end{proposition}

\begin{proof}
First, $\sigma(I_S)$ is dense in $G_S$: $S$ contains all primes up to $e^{n_0+1}$ by
Lemma~\ref{lem:bands}(1) and, by Lemma~\ref{lem:SW} applied to the principal character, meets
every reduced class modulo $M$ in every large band; so the classes of $S\setminus F$ generate
$(\Z/M\Z)^\times$ and $\chi\mapsto\tilde\chi$ is injective, whence $\tilde\chi$ is
non-principal. Now combine \eqref{eq:linear}, Lemma~\ref{lem:SW} and \eqref{eq:err}:
$\Lambda(x)=O(\sum_n|c_n|)+o(W(x))=o(W(x))$, so the left side of \eqref{eq:linear} is
$2W(x)(1+o(1))\to\infty$.
\end{proof}

\begin{proposition}[The trivial character]\label{prop:trivialchi}
$\sum_{\ell\in S}\ell^{-1}|1-\ell^{is}|^2$ is finite for $s\in2\pi\Z$ and infinite for
$s\notin2\pi\Z$.
\end{proposition}

\begin{proof}
Let $s=2\pi m$. For $\ell\in B_n$, $\log\ell=n+\delta_\ell$ with $n\in\Z$, so
$\ell^{is}=e^{2\pi im\delta_\ell}$ and
$|1-\ell^{is}|^2=4\sin^2(\pi m\delta_\ell)\le16\pi^2m^2\varepsilon_n^2$; hence the sum is
$\ll m^2\sum_n\varepsilon_n^2w_n\ll m^2\sum_n n^{-1}(\log n)^{-3/2}<\infty$.

Let $s\notin2\pi\Z$. Here $\tilde\chi$ is principal, so $c_n=w_n+O_M(1)\mathbf 1_{n\le n_F}$.
Since $w_n$ is monotone to $0$ up to an absolutely summable perturbation
(Lemma~\ref{lem:bands}(3)) and
$\bigl|\sum_{n\le N}e^{isn}\bigr|\le|\sin(s/2)|^{-1}<\infty$, Dirichlet's test gives
convergence of $\sum_ne^{isn}w_n$. With \eqref{eq:err}, $\Lambda(x)=O_s(1)+o(W(x))=o(W(x))$
and the sum diverges as before.
\end{proof}

\begin{lemma}[Matched pairs in consecutive bands]\label{lem:matched}
There are pairs $(a_j,b_j)$ of primes of $S$ with $a_j\in B_{n_j+1}$, $b_j\in B_{n_j}$, all
the $a_j,b_j$ distinct, such that $\log(a_j/b_j)\to1$, $\sigma_{b_j}\sigma_{a_j}^{-1}\to1$ in
$G_S=\hZ_S^\times$, and $\sum_ja_j^{-1}=\infty$.
\end{lemma}

\begin{proof}
For $\ell\in B_n$, $\log\ell=n+\delta_\ell$ with $|\delta_\ell|\le2\varepsilon_n$, so any
$a\in B_{n+1}$, $b\in B_n$ satisfy $|\log(a/b)-1|\le4\varepsilon_n\to0$. Put
$k_n=\lfloor\log\log n\rfloor$ and $M_n=\prod_{\ell'\le k_n}\ell'^{\,k_n}$, so that
$\log M_n\le2k_n^2=o(\log n)$. By Siegel--Walfisz, uniformly for moduli
$q\le(\log y)^{A}$ one has $\pi(y;q,a)=\mathrm{li}(y)/\varphi(q)+O_A(ye^{-c\sqrt{\log y}})$;
applied at the endpoints $e^{n\pm\varepsilon_n^\pm}$ of $B_n$ with $q=M_n\le n$, and since
$|B_n|\asymp\varepsilon_ne^n/n$ while the error is $O(e^{n-c\sqrt n})$, every reduced class
$a$ modulo $M_n$ contains at least $\tfrac12|B_n|/\varphi(M_n)$ primes of $B_n$ and at least
as many primes of $B_{n+1}$, for all large $n$. Pairing, injectively within each class,
primes of $B_{n+1}$ with primes of $B_n$ in the same class gives at least $\tfrac12|B_n|$
pairs $(a,b)\in B_{n+1}\times B_n$ with $a\equiv b\pmod{M_n}$; pairs from different $n$ are
disjoint because the bands are.

For a fixed prime $\ell'\in S$ and $k\ge1$ we have $\ell'^{\,k}\mid M_n$ for all large $n$,
and then the $\ell'$-component of $\sigma_b\sigma_a^{-1}$, which is $b/a$ for
$\ell'\notin\{a,b\}$ \cite[Lem.~2.10(1)]{Main}, is $\equiv1\pmod{\ell'^{\,k}}$; so every
coordinate of $\sigma_{b_j}\sigma_{a_j}^{-1}$ tends to $1$, which is convergence to $1$ in
the product topology of $\hZ_S^\times$. Finally, the pairs from band $n$ contribute
\[
\sum_{a}a^{-1}\ \ge\ \tfrac12|B_n|\,e^{-(n+2)}\ \ge\ \tfrac1{2e^{3}}\,w_n,
\qquad\text{since } w_n=\sum_{\ell\in B_n}\ell^{-1}\le|B_n|e^{-(n-1)},
\]
and $\sum_nw_n=\infty$ by Lemma~\ref{lem:bands}(4).
\end{proof}

\begin{proof}[Proof of Theorem~\ref{thm:counter}]
$\zeta_S(1)=\infty$ is Lemma~\ref{lem:bands}(4). $S$ is infinite, so
Lemma~\ref{lem:noexact} supplies two primes with irrationally related logarithms.
Propositions \ref{prop:chinontriv} and \ref{prop:trivialchi} give
$\Xi^{\mathrm{ext}}_1=\{1\}\times2\pi\Z$, hence $\Xi_1=\{1\}$ and $\Sigma_1=2\pi\Z$.
Proposition~\ref{prop:bound}(3) with $c=2\pi$ gives $E\subseteq\Z$: $M$ is of type
$\mathrm{III}_{e^{-k}}$ for some $k\ge1$ or of type $\mathrm{III}_0$, and in particular not
of type $\mathrm{III}_1$. Since $\Xi_1=\{1\}$, Lemma~\ref{lem:pairs} applies, and
Lemma~\ref{lem:matched} supplies pairs with $\beta\log(a_j/b_j)\to1$ ($\beta=1$),
$\sigma_{b_j}\sigma_{a_j}^{-1}\to1$ and $\sum a_j^{-1}=\infty$; hence $1\in E$, so $E=\Z$,
$S(M)=\{0\}\cup e^{\Z}$ by \eqref{eq:SM}, and $M$ is of type $\mathrm{III}_{e^{-1}}$.
\end{proof}

\begin{remark}[What the proof does not use]\label{rem:tools}
Only Mertens, Siegel--Walfisz at a \emph{fixed} modulus, Abel summation and Dirichlet's test
are needed. Because $\mathrm{round}(\log\ell)$ is constant on each band, the exponential sum
in \eqref{eq:linear} is over the \emph{integers} $n$, not over primes: its partial sums are
geometric and bounded by $|\sin(s/2)|^{-1}$, so Dirichlet's test handles all $s\notin2\pi\Z$
at once, with no distinction between rational and irrational $s/2\pi$ and no discrepancy
estimate. No Erd\H{o}s--Tur\'an inequality, no Vinogradov-type bound for
$\sum_\ell\ell^{-1+is}$, and no uniformity in the modulus is required. The bands are not
short intervals in the difficult sense: $B_n$ has multiplicative length $e^{2\varepsilon_n}$
with $\varepsilon_n\asymp(\log\log x)^{-1/2}$, so $|B_n|\gg x/(\log x)(\log\log x)^{1/2}$,
far above any threshold at which primes in progressions become hard.
\end{remark}

\begin{remark}[Numerical check]\label{rem:numtype}
The asymptotic regime of \eqref{eq:band} begins at $\varepsilon_n<\tfrac12$, i.e.\
$\ell>e^{e^4}\approx10^{23}$, and is beyond direct computation. The mechanics of
Lemmas~\ref{lem:bands}(3) and \ref{lem:SW} were instead verified with a fixed band
half-width $\varepsilon=0.05$, for which the same computations apply verbatim: over primes
up to $6\cdot10^7$ and bands $8\le n\le17$, the ratio
$w_n\big/\log\frac{n+\varepsilon}{n-\varepsilon}$ lies in $[0.998,1.003]$ for $n\ge13$, and
for the non-principal character modulo $4$ the ratio $|c_n|/w_n$ falls from
$8.6\cdot10^{-2}$ at $n=9$ to $2.1\cdot10^{-4}$ at $n=17$, consistent with
Lemma~\ref{lem:SW}.
\end{remark}

\begin{corollary}[Independence of the two classifications]\label{cor:indep}
Uniqueness of the $\mathrm{KMS}_\beta$ state and type $\mathrm{III}_1$ are logically
independent: Theorem~\ref{thm:counter} has the first and fails the second, while
Theorem~\ref{thm:classical} together with \cite[Cor.~3.10]{Main} has both. Equivalently,
$\Xiext$ can be trivial on the axis $s=0$ and nontrivial off it.
\end{corollary}

\section{The spectrum of transitions}\label{sec:spectrum}

We now turn to the horizontal direction: which chains $\{\Xi_\beta\}_{\beta>0}$ occur. By
\cite[Lem.~3.5]{Main} the chain is non-decreasing, and by \cite[Prop.~4.16]{Main} any
finite chain is realizable. We realize arbitrary countable chains.

\begin{theorem}[Realization]\label{thm:realize}
Let $(\beta_n)_{n\ge1}\subset(0,1)$ be any sequence. There is $S=Q\sqcup P_0\sqcup\bigsqcup_nR_n\subset P_\Q$
with
\begin{enumerate}
\item $\beta_c(S)=1$ and $\zeta_S(\beta)=\infty$ for $0<\beta\le1$;
\item $\sigma(I_S)$ dense in $G_S$;
\item with $q_n$ the $n$-th element of $Q$ and $\chi_n$ the quadratic character at that
component,
\[
\Xi_\beta(\Q,S)=\bigl\langle\chi_n:\beta_n<\beta\bigr\rangle\cong\bigoplus_{n:\ \beta_n<\beta}\Z/2\Z
\qquad(0<\beta\le1),
\]
and $\Xi_\beta=\widehat{G_S}$ for $\beta>1$. In particular the jump set is
$\{\beta_n\}\cup\{1\}$.
\end{enumerate}
\end{theorem}

The construction proceeds in blocks; the block scheme is what makes infinitely many
congruence conditions manageable.

\begin{lemma}[Independence of the quadratic conditions]\label{lem:disjoint}
Let $q_1<\dots<q_j$ be distinct odd primes, $q^*=(-1)^{(q-1)/2}q$, $M\ge1$, and
$E=\Q(\zeta_M)(\sqrt{q_1^*},\dots,\sqrt{q_j^*})$. Then
$[E:\Q]=\varphi(M)\cdot2^{\#\{i\le j:q_i\nmid M\}}$, and for any $a\in(\Z/M\Z)^\times$ and
signs $\epsilon_i$ compatible with $a$ at the $q_i$ dividing $M$, the primes $\ell$ with
$\ell\equiv a\ (M)$ and $\bigl(\tfrac{\ell}{q_i}\bigr)=\epsilon_i$ have natural density
$[E:\Q]^{-1}>0$.
\end{lemma}

\begin{proof}
The $q_i^*$ are multiplicatively independent modulo squares, so
$[\Q(\sqrt{q_1^*},\dots,\sqrt{q_j^*}):\Q]=2^j$ by Kummer theory. $\Q(\sqrt{q_i^*})$ is the
unique quadratic subfield of $\Q(\zeta_{q_i})$, so it lies in $\Q(\zeta_M)$ iff $q_i\mid M$;
the degree formula follows. Quadratic reciprocity turns the sign conditions into splitting
conditions, and Chebotarev applies.
\end{proof}

\begin{definition}[Block scheme]\label{def:blocks}
Choose $z_1<z_2<\cdots$ recursively, $z_{j+1}$ as large as needed in terms of $z_j$ and $j$,
and put $B_j=[z_j,z_{j+1})$. Let $q_j$ be a prime in $B_j$ with $\bigl(\tfrac{q_j}{q_i}\bigr)=+1$
for all $i<j$ (such primes have positive density by Lemma~\ref{lem:disjoint}, so they exist in
$B_j$ once $z_{j+1}$ is large), and $Q=\{q_j\}$. For $j\ge1$ set
\[
C_{0,j}=\bigl\{\ell:\ \bigl(\tfrac{\ell}{q_i}\bigr)=+1\ \forall i\le j\bigr\},
\quad
C_{n,j}=\bigl\{\ell:\ \bigl(\tfrac{\ell}{q_n}\bigr)=-1,\ \bigl(\tfrac{\ell}{q_i}\bigr)=+1\ \forall i\le j,\ i\ne n\bigr\},
\]
each of density $2^{-j}$ by Lemma~\ref{lem:disjoint}, and
\[
P_0=\bigsqcup_{j\ge1}\bigl(C_{0,j}\cap B_j\bigr)\setminus Q,
\qquad
R_n=\bigsqcup_{j\ge n}\rho_{n,j},\quad \rho_{n,j}\subseteq C_{n,j}\cap B_j\ \text{finite}.
\]
The essential property is that for $\ell\in S\cap[z_j,\infty)$ the symbols
$\bigl(\tfrac{\ell}{q_i}\bigr)$, $i\le j$, are determined by which block-class $\ell$ lies
in.
\end{definition}

\begin{lemma}[The pieces]\label{lem:pieces}
The $z_j$ and $\rho_{n,j}$ may be chosen so that:
\begin{enumerate}
\item $\sum_{\ell\in P_0\cap B_j}\ell^{-1}\ge1$ for every $j$, whence
$\sum_{P_0}\ell^{-1}=\infty$ and $\beta_c(S)=1$; and for every $M$ and every
$a\in(\Z/M\Z)^\times$ with $\bigl(\tfrac{a}{q_i}\bigr)=1$ whenever $q_i\mid M$,
$\sum_{\ell\in P_0,\ \ell\equiv a\,(M)}\ell^{-1}=\infty$;
\item $|R_n\cap[1,x]|\asymp_nx^{\beta_n}$, hence by Abel summation
$\sum_{\ell\in R_n}\ell^{-\beta}<\infty\iff\beta>\beta_n$, with divergence at the endpoint;
\item $q_n\ge z_n$ with $\log\log z_n\ge2^n$, so $\sum_nq_n^{-\beta}<\infty$ for every
$\beta>0$ and $Q$ is negligible in every convergence question.
\end{enumerate}
\end{lemma}

\begin{proof}
(1) $C_{0,j}$ has density $2^{-j}$, so
$\sum_{\ell\in C_{0,j},\ \ell\le x}\ell^{-1}=2^{-j}\log\log x+O_j(1)$; choose $z_{j+1}$ with
$2^{-j}(\log\log z_{j+1}-\log\log z_j)\ge2$. For the refined statement, for $j$ beyond the
largest index with $q_i\mid M$ the two conditions are compatible and
Lemma~\ref{lem:disjoint} gives positive density, so enlarging $z_{j+1}$ further (finitely
many extra requirements per block) gives block sums $\ge\varphi(M)^{-1}$.
(2) $C_{n,j}$ has positive density, so it meets $[y,2y]$ for $y\ge y_0(j)$; take
$\rho_{n,j}$ to contain one prime of $C_{n,j}$ in each $[m^{1/\beta_n},2m^{1/\beta_n}]$
contained in $B_j$. Abel summation with $A(x)=|R_n\cap[1,x]|\asymp x^{\beta_n}$ and $A(t)=0$ for $t<z_n$ gives
\[
\sum_{\ell\in R_n,\ \ell\le x}\ell^{-\beta}=A(x)x^{-\beta}+\beta\int_{z_n}^{x}A(t)t^{-\beta-1}\,dt
\ \asymp\ x^{\beta_n-\beta}+\beta\int_{z_n}^{x}t^{\beta_n-\beta-1}\,dt .
\]
The lower limit contributes the constant $z_n^{\beta_n-\beta}$; the integral converges as
$x\to\infty$ iff $\beta>\beta_n$, equals $\beta_n\log(x/z_n)\to\infty$ at $\beta=\beta_n$,
and diverges as a power for $\beta<\beta_n$.
(3) $q_n\ge\exp\exp(2^n)$ makes $q_n^{-\beta}\le\exp(-\beta\exp(2^n))$.
\end{proof}

\begin{proof}[Proof of Theorem~\ref{thm:realize}]
(1) is Lemma~\ref{lem:pieces}(1). For (2), fix a modulus $M$ supported on a finite
$F\subset S$, let $Q_F=Q\cap\{q:q\mid M\}$ and
$K_M=\{a:\bigl(\tfrac{a}{q}\bigr)=1\ \forall q\in Q_F\}$, a subgroup of index $2^{|Q_F|}$
with $(\Z/M\Z)^\times/K_M\cong(\Z/2\Z)^{Q_F}$. By Lemma~\ref{lem:pieces}(1) every $a\in K_M$
is realized by infinitely many $\ell\in P_0$, so $K_M\subseteq\langle\ell\bmod M\rangle$; and
for $q_n\in Q_F$ any large $\ell\in R_n$ outside $F$ maps to the basis vector $e_n$ of the
quotient. Hence the classes generate $(\Z/M\Z)^\times$.

For (3), let $\chi\ne1$ factor through $(\Z/M\Z)^\times$. If $\tilde\chi|_{K_M}\ne1$, pick
$a\in K_M$ with $\tilde\chi(a)\ne1$; by Lemma~\ref{lem:pieces}(1)
$\sum_{\ell\in P_0,\ \ell\equiv a}\ell^{-1}=\infty$ and every such $\ell$ detects $\chi$, so
$\chi\notin\Xi_\beta$ for $\beta\le1$. If $\tilde\chi|_{K_M}=1$ then
$\tilde\chi=\prod_{n\in T}\tilde\chi_{q_n}$ for a unique $T$, and by (2) and injectivity
$\chi=\chi_T:=\prod_{n\in T}\chi_n$. For $\ell\in S$ with $\ell\ge z_{\max T}$ the block
property gives $\chi_T(\sigma_\ell)=+1$ on $P_0$ and on $R_m$ with $m\notin T$, and $-1$ on
$R_m$ with $m\in T$; so, $Q$ being negligible and $[1,z_{\max T})$ finite,
\[
\sum_{\ell\in D_{\chi_T}}\ell^{-\beta}<\infty\iff\sum_{m\in T}\sum_{\ell\in R_m}\ell^{-\beta}<\infty
\iff\beta>\max_{m\in T}\beta_m,
\]
by Lemma~\ref{lem:pieces}(2). This is the stated formula; for $\beta>1$,
$\zeta_S(\beta)<\infty$ and \cite[Lem.~3.5]{Main} gives $\Xi_\beta=\widehat{G_S}$.
\end{proof}

\begin{remark}[Numerical check]\label{rem:numcascade}
Three ingredients were verified computationally on a scaled-down model with
$Q=\{7,11,13,17\}$ and primes up to $2\cdot10^6$: the $2^4=16$ joint sign patterns of
Lemma~\ref{lem:disjoint} occur with densities in $[0.0621,0.0629]$ against the predicted
$1/16=0.0625$; the detecting set of $\chi_T$ equals $\bigsqcup_{m\in T}R_m$ exactly, for
every subset $T$ tested; and a set thinned to counting function $\asymp x^{b}$ has its
$\beta$-weighted sums collapsing across $\beta=b$ as Lemma~\ref{lem:pieces}(2) requires.
\end{remark}

\section{Cantor transition loci, and the unattainable staircase}\label{sec:cantor}

\begin{definition}\label{def:locus}
The \emph{jump set} is $J(S)=\{\beta_\chi:\chi\ne1\}$, where $\beta_\chi$ is the convergence
exponent of the detecting set $D_\chi=\{\pp\in S:\chi(\sigma_\pp)\ne1\}$; the
\emph{transition locus} $\mathcal L(S)$ is the set of $\beta$ near which $\beta\mapsto\Xi_\beta$
is not constant.
\end{definition}

\begin{corollary}[Dense and Cantor loci]\label{cor:loci}
Taking $(\beta_n)$ an enumeration of $\Q\cap(0,1)$ in Theorem~\ref{thm:realize} gives
$J(S)=(\Q\cap(0,1))\cup\{1\}$ and $\mathcal L(S)=[0,1]$. Taking $(\beta_n)$ dense in a Cantor
set $\cC\subset(0,1)$ gives $\mathcal L(S)=\cC\cup\{1\}$: the transition locus is exactly the
prescribed Cantor set.
\end{corollary}

\begin{proof}
By Theorem~\ref{thm:realize}(3), $\Xi_\beta$ is constant on an open interval
$I\subseteq(0,1)$ if and only if $I$ contains no $\beta_n$. Hence a point of $(0,1)$ fails to
have a neighbourhood of constancy exactly when it is a limit of the $\beta_n$, so that
\[
\mathcal L(S)\cap(0,1)=\overline{J(S)\cap(0,1)}=\overline{\{\beta_n\}} ,
\]
identifying the non-local-constancy locus with the closure of the jump set, as in
Remark~\ref{rem:whatremains}. Finally $\beta=1$ always lies in $\mathcal L(S)$, since
$\Xi_\beta$ jumps to $\widehat{G_S}$ there.
\end{proof}

\begin{definition}\label{def:order}
For weights $w_n>0$ with $\sum w_n=1$ the \emph{normalized rank} is
$\dw_w(\beta)=\sum_{n:\beta_n<\beta}w_n\in[0,1]$.
\end{definition}

\begin{theorem}[Pure-jump staircase]\label{thm:staircase}
In the situation of Corollary~\ref{cor:loci} with locus $\cC$, $\dw_w$ is non-decreasing,
equals $0$ below $\min\cC$ and $1$ above $\max\cC$, is constant on every connected component
of $(0,1)\setminus\cC$, and its set of points of increase is exactly $\cC$.
\end{theorem}

\begin{proof}
A component of the complement contains no $\beta_n$, so $\dw_w$ is constant there; every
neighbourhood of a point of $\cC$ contains some $\beta_n$ by density.
\end{proof}

\begin{theorem}[The devil's staircase is unattainable]\label{thm:impossible}
For every set of primes $S$: $\widehat{G_S}$ is countable, hence $J(S)$ is countable; and
for every choice of weights $\dw_w$ is a pure jump function, discontinuous at every
$\beta_n$ with jump $w_n>0$, whose distributional derivative is the purely atomic measure
$\sum_nw_n\delta_{\beta_n}$. Consequently there is no $S$ and no weighting for which
$\dw_w$ equals the Cantor function, or any continuous non-constant function.
\end{theorem}

\begin{proof}
$G_S$ is a compact metrizable abelian group, being a Hausdorff quotient of the second
countable locally compact group $\mathbb I_S$ \cite[Prop.~2.5]{Main}; the dual of a compact
metrizable abelian group is countable. For each $\chi$ the set $\{\beta:\chi\in\Xi_\beta\}$ is an up-set, so each $\chi$
contributes a single jump and $J(S)=\{\beta_\chi\}$ is countable. In the situation of
Theorem~\ref{thm:realize}, $\chi_n\notin\Xi_{\beta_n}$ by the endpoint divergence in
Lemma~\ref{lem:pieces}(2) while $\chi_n\in\Xi_\beta$ for $\beta>\beta_n$, so the jump at
$\beta_n$ is at least $w_n>0$. A non-decreasing function of the form
$\sum_{\beta_n<\beta}w_n$ has purely atomic distributional derivative. The Cantor function
is continuous with non-atomic derivative, and a purely atomic probability measure and a
non-atomic one are mutually singular, hence distinct.
\end{proof}

\begin{theorem}[Approximation]\label{thm:approx}
Let $\cC$ be the middle-thirds Cantor set, $c$ the Cantor function, and $S$ as in
Corollary~\ref{cor:loci} with $(\beta_n)$ dense in $\cC$. For every $\varepsilon>0$ there
are weights $w_n>0$, $\sum w_n=1$, with $\|\dw_w-c\|_\infty<\varepsilon$.
\end{theorem}

\begin{proof}
Let $\mu_c$ be the Cantor measure. Partition $\cC$ into finitely many relatively clopen
pieces of diameter $<\delta$ and place the $\mu_c$-mass of each at a point $\beta_{n_i}$
inside it, giving $\nu=\sum_{i\le N}v_i\delta_{\beta_{n_i}}$. Since $c$ is continuous, weak
convergence of distribution functions is uniform (P\'olya's theorem, \cite[Problem~14.8]{Billingsley}), so for $\delta$ small
$\|F_\nu-c\|_\infty<\varepsilon/2$. Put $w_{n_i}=(1-\tfrac\varepsilon4)v_i$ and distribute
$\tfrac\varepsilon4$ among the remaining indices with arbitrary positive weights; then
$|\dw_w-F_\nu|\le\tfrac\varepsilon4F_\nu+\tfrac\varepsilon4\le\tfrac\varepsilon2$.
\end{proof}

\begin{remark}[Jump set versus transition locus]\label{rem:whatremains}
The obstruction in Theorem~\ref{thm:impossible} is not an artifact of the construction of
\S\ref{sec:spectrum}: $\Xi_\beta$ is a subgroup of the fixed countable group
$\widehat{G_S}$, and a monotone family of subgroups of a countable group can change at only
countably many parameters. What survives, and what Corollary~\ref{cor:loci} realizes, is the
distinction between the \emph{jump set}, necessarily countable, and the \emph{transition
locus} $\overline{J(S)}$, which is the closure of a countable set and hence an arbitrary
separable closed set --- in particular an arbitrary Cantor set. The Cantor spectrum of
transitions is real; the continuity of the staircase is not. Theorems
\ref{thm:impossible} and \ref{thm:approx} together say the devil's staircase is a limit
point of the realizable order parameters in the uniform norm, but not itself realizable.
\end{remark}

\section{Order parameters}\label{sec:order}

Each transition brings into existence a macroscopic observable, and the family is
independent.

\begin{definition}\label{def:Psin}
With $E_n=\{\ell\in S:\ell<z_n\}$ finite and
$E_n^-=\{\ell\in E_n:\bigl(\tfrac{\ell}{q_n}\bigr)=-1\}$, put
\begin{equation}\label{eq:Psin}
\Psi_n(x)=\Bigl(\tfrac{u_{q_n}(x)\bmod q_n}{q_n}\Bigr)\cdot(-1)^{\sum_{\ell\in R_n}v_\ell(x)}\cdot\prod_{\ell\in E_n^-}(-1)^{v_\ell(x)}\ \in\{\pm1\}.
\end{equation}
\end{definition}

\begin{proposition}\label{prop:Psin}
\begin{enumerate}
\item $\Psi_n(mx)=\Psi_n(x)$ for every $m\in\N_S$ and every $x$ in the domain --- pointwise,
not merely almost everywhere;
\item $\Psi_n$ is defined $\nu$-almost everywhere iff $\beta>\beta_n$;
\item under $\nusym$ the family $(\Psi_n)$ consists of independent uniform signs, and for
$\beta>\beta_n$ the map $\varphi\mapsto(\varphi(\Psi_n))_n$ is a bijection from the extreme
points of the $\mathrm{KMS}_\beta$ simplex onto $\{\pm1\}^{\{n:\beta_n<\beta\}}$.
\end{enumerate}
\end{proposition}

\begin{proof}
(1) It suffices to treat $m=\ell_0$ prime. Multiplication by $\ell_0$ sends
$v_\ell\mapsto v_\ell+\delta_{\ell,\ell_0}$ and $u_{q_n}\mapsto\ell_0u_{q_n}$ for
$\ell_0\ne q_n$ \cite[proof of Prop.~4.12]{Main}, so the first factor is multiplied by
$\bigl(\tfrac{\ell_0}{q_n}\bigr)$ unless $\ell_0=q_n$, when it is unchanged. If
$\ell_0=q_n$, no $v_\ell$ with $\ell\in R_n\cup E_n^-$ changes. If $\ell_0\in E_n$ and the
symbol is $-1$ then $\ell_0\in E_n^-$ and $v_{\ell_0}$ increases by one, flipping the third
factor. If $\ell_0\ge z_n$ and $\ell_0\notin R_n$, then either $\ell_0\in Q$, and the symbol
is $+1$ by the choice of $Q$ in Definition~\ref{def:blocks}, or the block property gives
symbol $+1$; in both cases no relevant valuation changes. If $\ell_0\in R_n$, the symbol is
$-1$ and $v_{\ell_0}$ increases by one, flipping the second factor.

(2) $\sum_{\ell\in R_n}\pi_\beta(v_\ell\ge1)=\sum_{\ell\in R_n}\ell^{-\beta}$ is finite iff
$\beta>\beta_n$ by Lemma~\ref{lem:pieces}(2); apply the two Borel--Cantelli lemmas to the
independent events $\{v_\ell\ge1\}$.

(3) Under $\nusym$ the conditional law of $u$ given $v$ is Haar, so the $u_{q_n}$ are
independent and each $\bigl(\tfrac{u_{q_n}}{q_n}\bigr)$ is a uniform sign independent of
everything else; multiplying by the remaining factors preserves uniformity and independence.
The bijection is Theorem~\ref{thm:realize}(3) together with \cite[Thm.~3.6]{Main}, which
identifies the extreme points with $\widehat{\Xi_\beta}=(\Z/2\Z)^{\{n:\beta_n<\beta\}}$.
\end{proof}

\begin{remark}[Numerical check]\label{rem:numpsi}
Take the model $Q=\{7,11,13,17\}$. Pointwise invariance was verified over
$4\cdot10^{5}$ random configurations, with no failures, together with a control
confirming that dropping the $R_n$-correction destroys it. The independence
asserted in (3) was checked by computing
$\mathbb E[\prod_{n\in T}\Psi_n]$ over all $15$ nonempty subsets of four indices, the largest
absolute value being $0.0037$ against a three-sigma bound of $0.0047$.
\end{remark}

\begin{remark}[Physical reading]\label{rem:memory}
Each $\Psi_n$ is a tail observable: changing the valuations at any finite set of primes
leaves it unchanged, because a change at $\ell\in R_n$ flips two factors of \eqref{eq:Psin}
simultaneously. It records one asymptotically conserved bit, coming into existence at the
sharp temperature $\beta_n$ where the Borel--Cantelli series for $R_n$ turns from divergent
to convergent. As $\beta$ increases across $\cC$ the bits condense on a countable dense
subset of a Cantor set, and $\dw_w(\beta)$ is the total weight of those already condensed;
between two condensations, that is on each complementary interval, the memory is frozen. By
Theorem~\ref{thm:impossible} they cannot condense continuously: each bit is a discrete
degree of freedom contributing a jump, and $\widehat{G_S}$ has only countably many. The
information-theoretic counterpart --- the relative entropy of an extremal state with respect
to the symmetric one grows by $\log2$ with each condensed bit --- is taken up in a companion
paper on invariants and entropy.
\end{remark}

\section{Assessment}

\[
\begin{array}{lll}
\text{Question} & \text{Answer} & \text{Reference}\\\hline
\text{factor type} & \text{bounded above by}\ \Sigma_\beta,\ \text{below by pairs} & \text{Prop.~\ref{prop:bound}, Lem.~\ref{lem:pairs}}\\
S=P_K,\ \beta\le1 & \mathrm{III}_1;\ \Sigma_\beta=\{0\}\ \text{via}\ L(1-i\tau,\chi)\ne0 & \text{Thm.~\ref{thm:classical}, Prop.~\ref{prop:slices}}\\
\text{two incommensurable norms} & \text{do not suffice} & \text{Thm.~\ref{thm:counter}}\\
\text{exact commensurability} & \text{impossible for infinite }S & \text{Lem.~\ref{lem:noexact}}\\
\text{uniqueness}\Rightarrow\mathrm{III}_1 & \text{false} & \text{Cor.~\ref{cor:indep}}\\
\text{chains}\ \{\Xi_\beta\} & \text{any}\ (\beta_n)\subset(0,1)\ \text{realizable} & \text{Thm.~\ref{thm:realize}}\\
\text{transition locus} & \text{any Cantor set} & \text{Cor.~\ref{cor:loci}}\\
\text{jump set} & \text{always countable} & \text{Thm.~\ref{thm:impossible}}\\
\text{devil's staircase} & \text{unattainable, approximable} & \text{Thms.~\ref{thm:impossible}, \ref{thm:approx}}
\end{array}
\]

Five questions remain. First, whether $\Sigma_\beta$ determines the type, i.e.\ whether the
flow of weights of these factors always has pure point spectrum (Remark~\ref{rem:flow}); if
it does, Proposition~\ref{prop:bound} is an equivalence and Lemma~\ref{lem:pairs} is not
needed. Second, the extended group $\Xiext$ has not been computed for the sparse sets
of \cite[\S5]{Main}; in particular whether Brun-summable $S$ gives $\mathrm{III}_\lambda$ or
$\mathrm{III}_0$ is open. In Theorem~\ref{thm:counter} we determined $\Sigma_1$ exactly, but
for a general band set with $\varepsilon(\ell)=(\log\log\ell)^{-\theta}$ the dependence of
$\Sigma_1$ on $\theta$ is unexamined. Fourth, Theorem~\ref{thm:impossible} shows the
realizable order parameters are pure jump functions; which monotone families of subgroups of
a countable group arise as $\{\Xi_\beta\}$, beyond the direct sums realized here, we do not
know. Finally, Lemma~\ref{lem:dict} says that $\Xi_\beta$ is the kernel of
$\chi\mapsto[\chi\circ\sigma]$ into the measurable cohomology group
$H^1_{\mathrm{meas}}(\cR_S,\pi_\beta;\T)$; that group itself is not computed here, and a
computation would presumably exhibit $\widehat{G_S}/\Xi_\beta$ inside it and give the chain
$\{\Xi_\beta\}$ a cohomological reading.

\begin{thebibliography}{9}
\bibitem{AW} H. Araki, E. J. Woods, \emph{A classification of factors}, Publ. RIMS Kyoto Univ. 4 (1968), 51--130.
\bibitem{Billingsley} P. Billingsley, \emph{Convergence of Probability Measures}, 2nd ed., Wiley, 1999.
\bibitem{BC} J.-B. Bost, A. Connes, \emph{Hecke algebras, type III factors and phase transitions with spontaneous symmetry breaking in number theory}, Selecta Math. (N.S.) 1 (1995), 411--457.
\bibitem{Connes73} A. Connes, \emph{Une classification des facteurs de type III}, Ann. Sci. \'Ec. Norm. Sup. (4) 6 (1973), 133--252.
\bibitem{CT} A. Connes, M. Takesaki, \emph{The flow of weights on factors of type III}, T\^ohoku Math. J. 29 (1977), 473--575.
\bibitem{Hecke} E. Hecke, \emph{Eine neue Art von Zetafunktionen und ihre Beziehungen zur Verteilung der Primzahlen}, Math. Z. 1 (1918), 357--376; 6 (1920), 11--51.
\bibitem{Krieger} W. Krieger, \emph{On ergodic flows and the isomorphism of factors}, Math. Ann. 223 (1976), 19--70.
\bibitem{Landau} E. Landau, \emph{\"Uber Ideale und Primideale in Idealklassen}, Math. Z. 2 (1918), 52--154.
\bibitem{LLN} M. Laca, N. S. Larsen, S. Neshveyev, \emph{On Bost--Connes type systems for number fields}, J. Number Theory 129 (2009), 325--338.
\bibitem{Maharam} D. Maharam, \emph{Incompressible transformations}, Fund. Math. 56 (1964), 35--50.
\bibitem{Neu} J. Neukirch, \emph{Algebraic Number Theory}, Grundlehren 322, Springer, 1999.
\bibitem{Schmidt} K. Schmidt, \emph{Cocycles on Ergodic Transformation Groups}, Macmillan Lectures in Math. 1, Delhi, 1977.
\bibitem{Main} R. Chen, \emph{KMS state uniqueness and phase transitions in localized Bost--Connes systems}, preprint, 2026, \newline\url{https://doi.org/10.5281/zenodo.22152101}.
\end{thebibliography}

\end{document}
